\begin{helper}
Fix a history cell \(C\), a finite terminal coordinate set \(U\), and two
disjoint fixed subsets \(P,Z\subseteq U\).  Suppose that, inside \(C\), every
complete terminal assignment \(A\subseteq U\) has conditional probability
\[
p^{|A|}(1-p)^{|U|-|A|}.
\]
If all coordinates of \(P\) are required to be present and all coordinates of
\(Z\) are required to be absent, then the conditional probability of this
fixed cylinder is at most
\[
p^{|P|}(1-p)^{|Z|}.
\]
\end{helper}
\begin{proof}
The event in question is the disjoint union of the complete terminal
assignment cells indexed by subsets \(A\subseteq U\) satisfying
\[
P\subseteq A,\qquad A\cap Z=\emptyset.
\]
For such an \(A\), write
\[
A=P\cup W,
\]
where
\[
W\subseteq U\setminus(P\cup Z).
\]
This representation is unique because \(P\) and \(Z\) are fixed and
disjoint.  The assumed complete-assignment product law gives the mass
\[
p^{|P|+|W|}(1-p)^{|Z|+|U|-|P|-|Z|-|W|}
\]
for the assignment corresponding to \(W\).  Factoring out the fixed
requirements gives
\[
p^{|P|}(1-p)^{|Z|}
\cdot
p^{|W|}(1-p)^{|U|-|P|-|Z|-|W|}.
\]
Summing over all subsets \(W\subseteq U\setminus(P\cup Z)\), the remaining
factor is
\[
\sum_{W\subseteq U\setminus(P\cup Z)}
p^{|W|}(1-p)^{|U|-|P|-|Z|-|W|}
=(p+(1-p))^{|U|-|P|-|Z|}=1.
\]
Thus the exact cylinder mass is
\[
p^{|P|}(1-p)^{|Z|}.
\]
If the conditioning cell has zero mass, the definition of conditional
probability makes the left side zero, so the same upper bound still holds.
\end{proof}
